\par\medskip
\hypertarget{native-069-article051}{}\pdfbookmark[2]{51.}{article-051}
\begin{center}51.\end{center}
Als Beispiel führen wir die in Art. 13 und 14 begonnene Rechnung weiter, zu der der Planet Juno die Zahlenwerte geliefert hatte. Oben hatten wir als wahre Anomalie $315^0\,1^{\prime}\,23^{\prime\prime}{,}02$ und als Logarithmus des Radiusvektors $0{,}325\,9877$ gefunden. Nun sei $i=13^0\,6^{\prime}\,44^{\prime\prime}{,}10$ und der Abstand des Perihels vom Knoten $=241^0\,10^{\prime}\,20^{\prime\prime}{,}57$, also $u=196^0\,11^{\prime}\,43^{\prime\prime}{,}59$. Schließlich sei $\AscendingNode=171^0\,7^{\prime}\,48^{\prime\prime}{,}73$. Daraus erhalten wir
\begin{center}
\begin{tabular}{@{}ll@{}}
\makebox[18em][l]{$\log\operatorname{tang}u$\hfill .\,.\,.\,.\,.\,.\,.} & $9{,}463\,0573$\\[.3em]
\makebox[18em][l]{$\log\cos i$\hfill .\,.\,.\,.\,.\,.\,.\,.} & $9{,}988\,5266$\\[.3em]
\hline
\makebox[18em][l]{$\log\operatorname{tang}(\lambda-\AscendingNode)$\hfill .\,.} & $9{,}451\,5839$\\[.3em]
$\lambda-\AscendingNode$\hfill $=$ & $195^0\,47^{\prime}\,40^{\prime\prime}{,}25$\\[.3em]
$\lambda=$ & $6\quad55\quad28{,}98$\\[.8em]
\makebox[18em][l]{$\log r$\hfill .\,.\,.\,.\,.\,.\,.\,.\,.\,.\,.} & $0{,}325\,9877$\\[.3em]
\makebox[18em][l]{$\log\cos\beta$\hfill .\,.\,.\,.\,.\,.\,.\,.} & $9{,}999\,1289$\\[.3em]
\hline
\makebox[18em][l]{$\log r^{\prime}$\hfill .\,.\,.\,.\,.\,.\,.\,.\,.\,.} & $0{,}325\,1166$
\end{tabular}
\end{center}
\begin{center}
\begin{tabular}{@{}ll@{}}
\makebox[18em][l]{$\log\sin(\lambda-\AscendingNode)$\hfill .\,.\,.\,.} & $9{,}434\,8692_n$\\[.3em]
\makebox[18em][l]{$\log\operatorname{tang}i$\hfill .\,.\,.\,.\,.\,.\,.\,.} & $9{,}367\,2304$\\[.3em]
\hline
\makebox[18em][l]{$\log\operatorname{tang}\beta$\hfill .\,.\,.\,.\,.\,.\,.} & $8{,}802\,0996_n$\\[.3em]
$\beta$\hfill $=$ & $-3^0\,37^{\prime}\,40^{\prime\prime}{,}02$\\[.3em]
\makebox[18em][l]{$\log\cos\beta$\hfill .\,.\,.\,.\,.\,.\,.\,.} & $9{,}999\,1289$\\[.3em]
\makebox[18em][l]{$\log\cos(\lambda-\AscendingNode)$\hfill .\,.\,.} & $9{,}983\,2852_n$\\[.3em]
\hline
 & $9{,}982\,4141_n$\\[.3em]
\makebox[18em][l]{$\log\cos u$\hfill .\,.\,.\,.\,.\,.\,.\,.} & $9{,}982\,4141_n$
\end{tabular}
\end{center}
\noindent Die Rechnung nach den Formeln III und VII würde sich folgendermaßen gestalten.
\begin{center}
\begin{tabular}{@{}ll@{}}
\makebox[18em][l]{$\log\sin u$\hfill .\,.\,.\,.\,.\,.\,.\,.\,.\,.\,.} & $9{,}445\,4714_n$\\[.3em]
\makebox[18em][l]{$\log\sin i$\hfill .\,.\,.\,.\,.\,.\,.\,.\,.\,.\,.} & $9{,}355\,7570$\\[.3em]
\hline
\makebox[18em][l]{$\log\sin\beta$\hfill .\,.\,.\,.\,.\,.\,.\,.\,.\,.\,.} & $8{,}801\,2284_n$\\[.3em]
$\beta=$ & $-3^0\,37^{\prime}\,40^{\prime\prime}{,}02$
\end{tabular}
\end{center}
\begin{center}
\begin{tabular}{@{}ll@{}}
\makebox[18em][l]{$\log\operatorname{tang}\tfrac{1}{2}i$\hfill .\,.\,.\,.\,.\,.\,.\,.\,.} & $9{,}060\,4258$\\[.3em]
\makebox[18em][l]{$\log\operatorname{tang}\beta$\hfill .\,.\,.\,.\,.\,.\,.\,.\,.\,.} & $8{,}802\,0996_n$\\[.3em]
\makebox[18em][l]{$\log\cos u$\hfill .\,.\,.\,.\,.\,.\,.\,.\,.\,.\,.} & $9{,}982\,4141_n$\\[.3em]
\hline
\makebox[18em][l]{$\log\sin(u-\lambda+\AscendingNode)$\hfill .\,.\,.} & $7{,}844\,9395$\\[.3em]
$u-\lambda+\AscendingNode=$ & $0^0\,24^{\prime}\,3^{\prime\prime}{,}34$\\[.3em]
$\lambda-\AscendingNode=$ & $195\quad47\quad40{,}25$
\end{tabular}
\end{center}

\SourceLeaf{070}\hypertarget{native-070-article052}{}\pdfbookmark[2]{52.}{article-052}
\begin{center}52.\end{center}
Betrachtet man $i$ und $u$ als veränderliche Größen, so erhält man durch Differenzieren der Gleichung III in Art. 50
\[
\operatorname{cotang}\beta\,\mathrm{d}\beta=\operatorname{cotang}i\,\mathrm{d}i+\operatorname{cotang}u\,\mathrm{d}u
\]
oder
\[
\mathrm{XII.}\quad\mathrm{d}\beta=\sin(\lambda-\AscendingNode)\,\mathrm{d}i+\sin i\cos(\lambda-\AscendingNode)\,\mathrm{d}u.
\]
Ebenso erhält man durch Differenzieren der Gleichung I
\[
\mathrm{XIII.}\quad\mathrm{d}(\lambda-\AscendingNode)=-\operatorname{tang}\beta\cos(\lambda-\AscendingNode)\,\mathrm{d}i+\frac{\cos i}{\cos\beta^2}\,\mathrm{d}u.
\]
Schließlich ergibt das Differenzieren der Gleichung XI
\[
\mathrm{d}r^{\prime}=\cos\beta\,\mathrm{d}r-r\sin\beta\,\mathrm{d}\beta
\]
oder
\[
\begin{aligned}
\mathrm{XIV.}\quad\mathrm{d}r^{\prime}={}&\cos\beta\,\mathrm{d}r-r\sin\beta\sin(\lambda-\AscendingNode)\,\mathrm{d}i\\
&{}-r\sin\beta\sin i\cos(\lambda-\AscendingNode)\,\mathrm{d}u.
\end{aligned}
\]
In dieser letzten Gleichung sind entweder die Terme mit $\mathrm{d}i$ beziehungsweise $\mathrm{d}u$ durch $206265^{\prime\prime}$ zu teilen oder die übrigen Terme mit dieser Zahl zu multiplizieren, wenn die Änderungen von $i$ und $u$ in Bogensekunden angegeben werden.
